\section*{Problemas}

\subsection*{Enunciados de los problemas}

% Recordar
\begin{problema}
	\label{prob:054ed26}
	Encontrar la probabilidad de que un solo lanzamiento de un dado balanceado de 6 caras resulte en un número menor que 4 si:
	\begin{enumerate}
		\item no hay más información disponible;
		\item se sabe con certeza que el lanzamiento resultó en un número impar.
	\end{enumerate}
\end{problema}

% Comprender
\begin{problema}
	\label{prob:bde7784}
	En un grupo de estudiantes de estadística, el 60\% son mujeres ($M$) y el 40\% son hombres ($H$). Se sabe que el 10\% de las mujeres son zurdas ($Z$) y el 15\% de los hombres son zurdos. Si se selecciona una persona al azar y resulta ser zurda, calcula aplicando el Teorema de Bayes la probabilidad de que sea mujer, y explica por qué este resultado no es simplemente igual a la proporción original de mujeres en el grupo ($60\%$).
\end{problema}

% Aplicar
\begin{problema}
	\label{prob:e14faa0}
	Un estudiante se prepara para un examen general. Si estudia debidamente ($E$), la probabilidad de aprobar ($A$) es de 0.90. Si no estudia ($E'$), la probabilidad de aprobar es de sólo 0.30. La probabilidad de que el estudiante estudie es de 0.70.
	\begin{enumerate}
		\item ¿Cuál es la probabilidad total de que el estudiante apruebe el examen?
		\item Si el estudiante aprobó, ¿cuál es la probabilidad condicional de que haya estudiado?
		\item Si el estudiante no aprobó, ¿cuál es la probabilidad de que no haya estudiado?
	\end{enumerate}
\end{problema}

% Analizar
\begin{problema}
	\label{prob:a7e87e4}
	Demuestra rigurosamente el Teorema de Bayes para una partición finita y exhaustiva $\{B_1, B_2, \dots, B_k\}$ del espacio muestral $S$ dado un evento observable $A$ con $P(A)>0$:
	\begin{align}
		P(B_j | A) = \frac{P(A|B_j)P(B_j)}{\sum_{i=1}^k P(A|B_i)P(B_i)}, \quad \text{para todo } j \in \{1, \dots, k\}.
	\end{align}
\end{problema}

% Evaluar
\begin{problema}
	\label{prob:898898e}
	\textbf{La Falacia del Fiscal y la razón de momios:} En una investigación forense se recupera una muestra biológica en la escena del crimen. El perfil de ADN extraído tiene una probabilidad de coincidencia aleatoria en la población de $\alpha = 10^{-6}$ (es decir, $P(\text{coincidencia} | \text{inocente}) = 10^{-6}$). Un sospechoso es identificado mediante un rastreo masivo y su ADN coincide perfectamente con la muestra ($P(\text{coincidencia} | \text{culpable}) = 1$). El fiscal afirma ante el jurado que, puesto que la probabilidad de coincidencia aleatoria es de apenas una en un millón, la probabilidad de que el sospechoso sea inocente es del $0.0001\%$.

	Evalúa mediante la formulación en momios (odds ratio) del Teorema de Bayes por qué el razonamiento del fiscal es erróneo. Si en la región habitan $N = 5,000,000$ de personas con posible acceso a la escena y no existe ninguna otra evidencia contra el sospechoso, calcula la verdadera probabilidad a posteriori de culpabilidad $P(\text{culpable} | \text{coincidencia})$.
\end{problema}

% Crear
\begin{problema}
	\label{prob:3bf9f42}
	Diseñe un problema original de clasificación de riesgo crediticio: un banco clasifica a sus solicitantes de préstamo en tres categorías de historial —Excelente, Regular y Deficiente— que representan el $50\%$, $35\%$ y $15\%$ de los solicitantes, respectivamente. Las tasas históricas de incumplimiento (\emph{default}) son del $2\%$, $8\%$ y $25\%$ en cada categoría. Calcule el vector completo de probabilidades a posteriori dado que un préstamo entró en incumplimiento, e identifique cuál categoría es la fuente más probable de incumplimientos en términos absolutos, a pesar de no ser la más numerosa.
\end{problema}

\subsection*{Soluciones detalladas}

% Recordar
\begin{solproblema}[prob:054ed26]
	\begin{enumerate}
		\item \textbf{Sin información:} el espacio muestral completo tiene 6 eventos equiprobables $S=\{1,2,3,4,5,6\}$. Los resultados estrictamente menores que 4 son $A=\{1,2,3\}$, por lo que $P(A) = \frac{3}{6} = 0.50$.
		\item \textbf{Con información de número impar:} el espacio condicionado es $I=\{1,3,5\}$. Los elementos de $A$ que también están en $I$ son $A \cap I = \{1,3\}$. Aplicando la definición de probabilidad condicional:
		\begin{align}
			P(A|I) = \frac{P(A \cap I)}{P(I)} = \frac{2/6}{3/6} = \frac{2}{3} \approx 0.6667.
		\end{align}
	\end{enumerate}
\end{solproblema}

% Comprender
\begin{solproblema}[prob:bde7784]
	Sean los eventos $M$ (mujer) con $P(M)=0.60$ y $H$ (hombre) con $P(H)=0.40$. Las tasas de zurdería son $P(Z|M)=0.10$ y $P(Z|H)=0.15$.

	Por la regla de la probabilidad total, la probabilidad de que una persona elegida al azar sea zurda es:
	\begin{align}
		P(Z) = P(Z|M)P(M) + P(Z|H)P(H) = (0.10)(0.60) + (0.15)(0.40) = 0.060 + 0.060 = 0.120.
	\end{align}
	Por el Teorema de Bayes:
	\begin{align}
		P(M|Z) = \frac{P(Z|M)P(M)}{P(Z)} = \frac{0.060}{0.120} = 0.50 \text{ (50\%).}
	\end{align}
	Este resultado ($50\%$) es menor que la proporción original de mujeres en el grupo ($60\%$) porque \emph{observar que la persona es zurda} es evidencia que favorece relativamente más a los hombres: aunque hay más mujeres en total, los hombres tienen una tasa de zurdería más alta ($15\%$ contra $10\%$), así que condicionar sobre "ser zurdo" desplaza la probabilidad hacia el grupo con mayor propensión relativa a ese rasgo, no solo hacia el grupo más numeroso.
\end{solproblema}

% Aplicar
\begin{solproblema}[prob:e14faa0]
	Sea $A$ el evento de aprobar y $E$ el de estudiar ($P(E)=0.70$, $P(E')=0.30$). Las probabilidades condicionales del examen son $P(A|E)=0.90$ y $P(A|E')=0.30$.
	\begin{enumerate}
		\item Probabilidad marginal total de aprobar:
		\begin{align}
			P(A) = P(A|E)P(E) + P(A|E')P(E') = (0.90)(0.70) + (0.30)(0.30) = 0.63 + 0.09 = 0.72 \text{ (72\%).}
		\end{align}
		\item Por el Teorema de Bayes, si aprobó, la probabilidad de haber estudiado es:
		\begin{align}
			P(E|A) = \frac{P(A|E)P(E)}{P(A)} = \frac{0.63}{0.72} = 0.875 \text{ (87.5\%).}
		\end{align}
		\item Para $A'$: $P(A') = 1-0.72=0.28$; $P(A'|E)=0.10$, $P(A'|E')=0.70$. Por Bayes sobre los complementos:
		\begin{align}
			P(E'|A') = \frac{P(A'|E')P(E')}{P(A')} = \frac{(0.70)(0.30)}{0.28} = \frac{0.21}{0.28} = 0.75 \text{ (75\%).}
		\end{align}
	\end{enumerate}
\end{solproblema}

% Analizar
\begin{solproblema}[prob:a7e87e4]
	Sea $\{B_1, B_2, \dots, B_k\}$ una partición exhaustiva ($\bigcup_{i=1}^k B_i = S$) y mutuamente excluyente. Por la regla del producto de la probabilidad condicional, para cualquier índice $j$:
	\begin{align}
		P(B_j \cap A) = P(B_j | A) P(A) = P(A | B_j) P(B_j).
	\end{align}
	Al despejar $P(B_j|A)$ dado que $P(A)>0$:
	\begin{align}
		P(B_j | A) = \frac{P(A | B_j) P(B_j)}{P(A)}. \label{eq:bayes_num_tdb}
	\end{align}
	Como los eventos $B_i$ forman una partición de $S$, expresamos $A$ como la unión disjunta de sus intersecciones sobre cada celda:
	\begin{align}
		A = A \cap S = A \cap \left( \bigcup_{i=1}^k B_i \right) = \bigcup_{i=1}^k (A \cap B_i).
	\end{align}
	Por el tercer axioma de Kolmogorov (aditividad para eventos disjuntos):
	\begin{align}
		P(A) = \sum_{i=1}^k P(A \cap B_i) = \sum_{i=1}^k P(A | B_i) P(B_i). \label{eq:prob_total_k_tdb}
	\end{align}
	Sustituyendo (\ref{eq:prob_total_k_tdb}) en (\ref{eq:bayes_num_tdb}), se obtiene la fórmula general del Teorema de Bayes para $k$ estados.
\end{solproblema}

% Evaluar
\begin{solproblema}[prob:898898e]
	Sea $I$ el evento de ser el culpable real del crimen, y $E$ la coincidencia en el perfil de ADN. El fiscal incurre en una falacia condicional al igualar erróneamente $P(E|I') = 10^{-6}$ con la probabilidad inversa $P(I'|E)$.

	Evaluamos la veracidad formal mediante la versión del Teorema de Bayes en razón de momios (Odds Form):
	\begin{align}
		\frac{P(I|E)}{P(I'|E)} = \underbrace{\frac{P(E|I)}{P(E|I')}}_{\text{Likelihood Ratio (LR)}} \times \underbrace{\frac{P(I)}{P(I')}}_{\text{Momios a priori}}.
	\end{align}
	Sin más evidencias que la ubicación geográfica en la población de $N = 5,000,000$ habitantes, la probabilidad a priori de que un individuo arbitrario sea culpable es $P(I) = \frac{1}{5,000,000} = 2 \times 10^{-7}$, con $P(I') \approx 1$. El cociente de verosimilitud es:
	\begin{align}
		\text{LR} = \frac{P(E|I)}{P(E|I')} = \frac{1}{10^{-6}} = 1,000,000.
	\end{align}
	Los momios a posteriori son:
	\begin{align}
		\text{Momios}(I|E) = 1,000,000 \times \frac{1/5,000,000}{1} = \frac{1}{5} = 0.20.
	\end{align}
	Convirtiendo a probabilidad exacta:
	\begin{align}
		P(I|E) = \frac{0.20}{1+0.20} = \frac{1}{6} \approx 0.1667 \text{ (16.67\%).}
	\end{align}
	La probabilidad de que el sospechoso sea culpable es de sólo $16.67\%$ (no del $99.9999\%$ como alegaba el fiscal). El error radica en no considerar que, al evaluar un perfil en una población de 5 millones donde la frecuencia de coincidencia aleatoria es 1 en 1 millón, se esperan estadísticamente unos 5 falsos positivos puramente aleatorios en la población —el sospechoso es apenas uno entre esos 5 candidatos igualmente compatibles, más el verdadero culpable.
\end{solproblema}

% Crear
\begin{solproblema}[prob:3bf9f42]
	Sean las categorías $Exc$, $Reg$, $Def$ con $P(Exc)=0.50$, $P(Reg)=0.35$, $P(Def)=0.15$, y tasas de incumplimiento $P(D|Exc)=0.02$, $P(D|Reg)=0.08$, $P(D|Def)=0.25$.

	La probabilidad marginal de incumplimiento es:
	\begin{align}
		P(D) = (0.02)(0.50)+(0.08)(0.35)+(0.25)(0.15) = 0.010+0.028+0.0375 = 0.0755.
	\end{align}
	El vector de probabilidades a posteriori es:
	\begin{align}
		P(Exc|D) &= \frac{0.010}{0.0755} \approx 0.1325 \text{ (13.25\%)} \\
		P(Reg|D) &= \frac{0.028}{0.0755} \approx 0.3709 \text{ (37.09\%)} \\
		P(Def|D) &= \frac{0.0375}{0.0755} \approx 0.4967 \text{ (49.67\%)}
	\end{align}
	La categoría \textbf{Deficiente} es la fuente más probable de incumplimientos en términos absolutos ($49.67\%$), a pesar de representar apenas el $15\%$ de los solicitantes: su tasa de incumplimiento ($25\%$) es tan desproporcionadamente alta que domina el cálculo, de forma análoga a como una línea de producción pequeña pero muy defectuosa puede superar en piezas defectuosas absolutas a una línea grande pero confiable.
\end{solproblema}
